%ID VD_4_1_2_01
\textbf{Bước 1. Tách biến theo hai biên Dirichlet.}
Với $u_t=Ku_{xx}$, $u(0,t)=u(L,t)=0$, đặt $u=XT$. Ta có $X''+\lambda X=0$, $T'+K\lambda T=0$. Điều kiện $X(0)=X(L)=0$ cho $X_n=\sin(n\pi x/L)$, $\lambda_n=(n\pi/L)^2$, chuẩn $L/2$. Vì vậy $u=\sum c_ne^{-K\lambda_nt}X_n$, $c_n=(2/L)\int_0^LfX_n$.

\textbf{Bước 2. Tính riêng từng dữ liệu.}
(i) $K=L=1$. Từ $\sin3z=3\sin z-4\sin^3z$ suy $\sin^3(\pi x)=3\sin(\pi x)/4-\sin(3\pi x)/4$. Trực giao cho $c_1=3/4,c_3=-1/4$, các hệ số khác bằng không. Do đó
\[u=\frac34e^{-\pi^2t}\sin\pi x-\frac14e^{-9\pi^2t}\sin3\pi x.\]
Tại $t=0$, công thức trả lại $f$; hai sine bằng không ở hai đầu. Đạo hàm mỗi hạng theo $t$ bằng đạo hàm theo $x$ cấp hai, nên PDE đúng. Đây là nghiệm trơn đến $t=0$.

(ii) $K=2,L=\pi$, nhưng $f=\sin(\pi x)$, không phải $\sin x$. Dùng $2\sin(\pi x)\sin(nx)=\cos((\pi-n)x)-\cos((\pi+n)x)$:
\begin{align*}
c_n&=\frac1\pi\left[\frac{\sin((\pi-n)x)}{\pi-n}-\frac{\sin((\pi+n)x)}{\pi+n}\right]_0^\pi\\
&=\frac{(-1)^n\sin(\pi^2)}{\pi}\left(\frac1{\pi-n}-\frac1{\pi+n}\right)
=\frac{2(-1)^nn\sin(\pi^2)}{\pi(\pi^2-n^2)}.
\end{align*}
Không có mẫu bằng không vì $\pi$ không nguyên. Nghiệm là
\[u=\frac{2\sin(\pi^2)}{\pi}\sum_{n=1}^\infty\frac{(-1)^nn}{\pi^2-n^2}e^{-2n^2t}\sin nx.\]
\textbf{Bước 3. Kiểm tra ý nghĩa nghiệm.} Chuỗi trơn với $t>0$, có biên $0$, và tiến về $f(x)$ khi $0<x<\pi$, $t\downarrow0$. Tuy nhiên $f(\pi)=\sin(\pi^2)\ne0$: không có nghiệm liên tục đến góc $(\pi,0)$ thỏa mọi dữ kiện. Khái niệm nghiệm thông thường nếu chỉ yêu cầu liên tục trên toàn bao đóng không được thỏa ở câu này.

%ID VD_4_1_2_02
\textbf{Bước 1. Cơ sở Neumann.} $X'(0)=X'(L)=0$ cho mode hằng $1$ và $X_n=\cos(n\pi x/L)$, $n\ge1$. Chuẩn lần lượt $L,L/2$, nên
\[u=c_0+\sum_{n\ge1}c_ne^{-K(n\pi/L)^2t}\cos(n\pi x/L),\quad c_0=\frac1L\int f,\quad c_n=\frac2L\int f\cos(n\pi x/L).\]
\textbf{Bước 2. Dữ liệu (i).} $\cos^2(\pi x)=(1+\cos2\pi x)/2$, nên với $K=L=1$, $c_0=c_2=1/2$, các hệ số khác bằng không:
\[u=\frac12+\frac12e^{-4\pi^2t}\cos2\pi x.\]
Đạo hàm biên là $-\pi e^{-4\pi^2t}\sin2\pi x$, bằng không ở $0,1$. Thế $t=0$ trả lại $\cos^2\pi x$; đạo hàm $u_t,u_{xx}$ đều bằng $-2\pi^2e^{-4\pi^2t}\cos2\pi x$. Nghiệm trơn và tiến đều về $1/2$.

\textbf{Bước 3. Dữ liệu (ii).} $K=2,L=\pi,f=\sin x$, $c_0=2/\pi$. Với $n\ne1$,
\begin{align*}
\int_0^\pi\sin x\cos nx\,dx
&=\frac12\int_0^\pi[\sin((1+n)x)+\sin((1-n)x)]dx\\
&=\frac{1+(-1)^n}{1-n^2}.
\end{align*}
Khi $n=1$, $\int\sin x\cos x=[\sin^2x/2]_0^\pi=0$. Vì thế $c_n=0$ khi $n$ lẻ, $c_{2j}=4/[\pi(1-4j^2)]$. Suy
\[u=\frac2\pi+\frac4\pi\sum_{j\ge1}\frac{e^{-8j^2t}}{1-4j^2}\cos2jx.\]
\textbf{Bước 4. Kiểm tra tính thông thường.} Hệ số $O(j^{-2})$ cho hội tụ đều đến $t=0$, trả lại $f$ liên tục. Khi $t>0$, chuỗi đạo hàm hội tụ đều, đúng PDE và Neumann. Vì $f'(0)=1,f'(\pi)=-1$, đạo hàm không thể liên tục đến hai góc và đồng thời bằng $0$ trên biên. Nghiệm vẫn thuộc lớp $C^{2,1}$ với $t>0$ và liên tục trên bao đóng, là lớp nghiệm thông thường được dùng trong ví dụ; không khẳng định trơn đến góc. Ước lượng
\[\left|u-\frac2\pi\right|\le\frac4\pi e^{-8t}\sum_{j\ge1}\frac1{4j^2-1}\longrightarrow0\]
chứng minh giới hạn đều $2/\pi$.

%ID VD_4_1_2_03
\textbf{Bước 1. Năng lượng và duy nhất.} Biên là $u_x(0)=u(0)$, $u_x(L)=0$. Tích phân từng phần:
\[I'=\int uu_t=K[uu_x]_0^L-K\int u_x^2
=-K u(0,t)^2-K\int_0^Lu_x^2dx\le0.\]
Áp dụng cho hiệu hai nghiệm có dữ liệu đầu $0$ cho $I(0)=0$, nên hiệu bằng không.

\textbf{Bước 2. Trị riêng và trực giao.} Thế $u=XT$ cho $X''+\lambda X=0$, $X'(0)=X(0)$, $X'(L)=0$, $T'=-K\lambda T$. Nhân ODE với $X$:
\[\lambda\int X^2=\int(X')^2+X(0)^2.\]
Nếu $\lambda=0$, $X'=0$ và $X(0)=0$, suy $X=0$; mọi trị riêng không tầm thường có $\lambda>0$. Trừ hai phương trình như đẳng thức Green cho $(\lambda_n-\lambda_m)\int X_nX_m=[X_nX_m'-X_n'X_m]_0^L=0$, nên trực giao.

\textbf{Bước 3. Tìm $k$ và chuẩn hóa.} Với $\lambda=k^2$, $X=A\cos kx+B\sin kx$. Biên trái cho $A=kB$; chọn $A=1$, được $X=\cos kx+\sin kx/k$. Biên phải cho $-k\sin kL+\cos kL=0$, tức $k\tan kL=1$. Trên mỗi khoảng $((n-1)\pi/L,(n-1/2)\pi/L)$, $k\tan kL$ tăng từ $0$ đến $\infty$, nên có đúng một nghiệm $k_n$, $n\ge1$. Chuẩn là
\[N(k)=\frac L2(1+k^{-2})+\frac{\sin2kL}{4k}(1-k^{-2})+\frac{\sin^2kL}{k^2}.\]
Do đó $u=\sum c_ne^{-Kk_n^2t}X_n$, $c_n=\int fX_n/N(k_n)$.

\textbf{Bước 4. Tính dữ liệu cụ thể $K=2,L=\pi,f=1+2\sin(x/2)$.}
Đặt $q=1/2$, và với $r\ne0$ đặt $D(r)=(1-\cos\pi r)/r$, $E(r)=\sin\pi r/r$, mở rộng $D(0)=0,E(0)=\pi$ theo giới hạn. Tích–tổng và tích phân từng hạng cho
\begin{align*}
J(k)=\int_0^\pi fX\,dx
={}&\frac{\sin\pi k}{k}+\frac{1-\cos\pi k}{k^2}\\
&+D(q+k)+D(q-k)+\frac{E(q-k)-E(q+k)}k.
\end{align*}
Vậy $c_n=J(k_n)/N(k_n)$ và
\[u(x,t)=\sum_{n\ge1}\frac{J(k_n)}{N(k_n)}e^{-2k_n^2t}
\left(\cos k_nx+\frac{\sin k_nx}{k_n}\right).\]
Không còn tích phân hệ số chưa tính. Với $L=\pi$, các nghiệm đầu là $0.3834486028,1.2187220137,2.1391968186$. Mỗi hạng thỏa hai biên và PDE. $f'(0)=1=f(0)$, $f'(\pi)=0$, nên dữ liệu khớp biên; chuỗi trả lại $f$ theo khai triển Sturm–Liouville và trơn khi $t>0$. Tất cả $k_n>0$, suy giới hạn đều bằng $0$. Không mặc định các đạo hàm cấp cao liên tục ở góc, vì điều đó cần thêm tương thích.

%ID VD_4_2_2_02
\textbf{Bước 1. Cận dữ liệu và rào trên đoạn lớn.} Giả thiết $u$ bị chặn trên $\mathbb R\times[0,T]$. Đặt $M=\sup_{\mathbb R\times[0,T]}u$, $N=\sup_xu(x,0)$ và cố định $(x_0,t_0)$, $0<t_0\le T$. Với $\varepsilon>0$,
\[w_\varepsilon=u-\varepsilon\{(x-x_0)^2+2Kt\}.\]
Tính $(w_\varepsilon)_t-K(w_\varepsilon)_{xx}=0$, vì $2K-K\cdot2=0$.
\textbf{Bước 2. Kiểm tra toàn bộ biên.} Trên mặt $t=0$, $w_\varepsilon\le N$. Trên $x=x_0\pm R$, $w_\varepsilon\le M-\varepsilon R^2$. Chọn $R>|x_0|+1$ đủ lớn để $M-\varepsilon R^2\le N$; áp dụng cực đại trên $(x_0-R,x_0+R)\times(0,t_0]$ được $w_\varepsilon\le N$.
\textbf{Bước 3. Bỏ rào và chứng minh cận dưới.} Tại điểm cố định,
\[u(x_0,t_0)\le N+2K\varepsilon t_0.\]
Cho $\varepsilon\downarrow0$, suy $u\le N$. Áp dụng cùng lập luận cho $-u$ cho $u\ge\inf_xu(x,0)$. Do $(x_0,t_0)$ tùy ý,
\[\inf_yu(y,0)\le u(x,t)\le\sup_yu(y,0).\]
Giả thiết bị chặn chính là điều cho phép chọn $R$ bằng cận $M$ độc lập theo $x$; không bỏ bước này trên miền vô hạn.

%ID VD_4_4_2_01
\textbf{Bước 1. Công thức mode có nguồn.}
Với Dirichlet thuần nhất trên $(0,\pi)$, đặt $u=\sum T_n\sin nx$. Chiếu PDE cho $T_n'+n^2T_n=F_n$, $T_n(0)=f_n$. Nhân với $e^{n^2t}$ rồi tích phân:
\[T_n=f_ne^{-n^2t}+\int_0^te^{-n^2(t-s)}F_n(s)ds.\]
\textbf{Bước 2. Ví dụ 1(i), $F=1,f=0$.}
\[F_n=\frac2\pi[-\cos ny/n]_0^\pi=\frac{2(1-(-1)^n)}{\pi n}.\]
Vì $\int_0^te^{-n^2(t-s)}ds=(1-e^{-n^2t})/n^2$,
\[u=\sum_{n\ge1}\frac{2(1-(-1)^n)}{\pi n^3}(1-e^{-n^2t})\sin nx.\]
Trạng thái dừng giải $P''=-1$, $P(0)=P(\pi)=0$, nên $P=x(\pi-x)/2$. Hệ số sine của $P$ chính là $F_n/n^2$; do đó $u\to P$, không tiến về hằng $1$. Đồ thị trong sách cần được đối chiếu với nghiệm này.

\textbf{Bước 3. Ví dụ 1(ii), biên $T_1,T_2$, nguồn và dữ liệu đầu bằng không.}
Đặt $b=T_1+(T_2-T_1)x/\pi$ và $w=u-b$. $b_t=b_{xx}=0$, nên $w_t=w_{xx}$, biên $0$, dữ liệu đầu $-b$. Tích phân từng phần:
\[\int_0^\pi b\sin nx\,dx
=\left[-\frac{b\cos nx}n\right]_0^\pi+\frac{b'}n\int_0^\pi\cos nx\,dx
=\frac{T_1-(-1)^nT_2}n.\]
Vì vậy
\[u=b-\sum_{n\ge1}\frac{2(T_1-(-1)^nT_2)}{\pi n}e^{-n^2t}\sin nx.\]
Biên được bảo đảm bởi $b$, mặt đầu trong miền bằng $0$ theo khai triển sine của $b$, và $u\to b$ đều lúc thời gian dài. Nếu $T_1$ hoặc $T_2$ khác $0$, không có liên tục tại góc tương ứng với dữ liệu đầu $0$.

\textbf{Bước 4. Ví dụ 2 nằm tiếp trong cùng phần scan: Neumann có nguồn.}
Với $u_x(0)=\mu,u_x(L)=\nu$, chọn $b=\mu x+(\nu-\mu)x^2/(2L)$. Khi đó $w=u-b$ có biên Neumann $0$, nguồn $Q=F+K(\nu-\mu)/L-b_t$, dữ liệu $f-b_0$. Nhân Green Neumann cho $w=\int G_N(f-b_0)+\int_0^t\int G_NQ$. Cũng có công thức trực tiếp
\[u=\int_0^LG_N(x,y,t)f(y)dy+
\int_0^t\!\int_0^LG_N(x,y,t-s)F(y,s)dy\,ds+
K\int_0^t\{G_N(x,L,t-s)\nu(s)-G_N(x,0,t-s)\mu(s)\}ds.\]
Dấu hai đầu được suy từ $K[G_Nu_y-(G_N)_yu]_0^L$.
Với $K=2,L=1,\mu=\nu=f=0,F=e^{-t}\cos^2\pi x$, dùng $\cos^2\pi x=(1+\cos2\pi x)/2$. Chỉ có mode $0,2$:
\[T_0'=\tfrac12e^{-t},\quad T_2'+8\pi^2T_2=\tfrac12e^{-t},\quad T_0(0)=T_2(0)=0.\]
Tích phân cho
\[u=\frac{1-e^{-t}}2+\frac{e^{-t}-e^{-8\pi^2t}}{2(8\pi^2-1)}\cos2\pi x.\]
Đúng biên, nguồn, mặt đầu; giới hạn bằng $1/2$. Hai ví dụ có đánh số 1 và 2 được giữ chung ở đây vì cùng nằm trong ảnh đề của mục này.

%ID VD_4_5_2_01
\textbf{Bước 1. Tính phép đổi ẩn.}
Đặt $u=e^{-hx}v$, $h=A/(2K)$. Khi đó $u_x=e^{-hx}(v_x-hv)$, $u_{xx}=e^{-hx}(v_{xx}-2hv_x+h^2v)$, $u_t=e^{-hx}v_t$. Thay vào $u_t=Ku_{xx}+Au_x+Bu$ cho
\[v_t=Kv_{xx}+(B-A^2/(4K))v,\quad v(x,0)=e^{hx}f(x).\]
Hai biên Dirichlet vẫn bằng $0$, nên $u=e^{-hx}\sum c_ne^{[B-A^2/(4K)-Kn^2]t}\sin nx$ khi $L=\pi$.

\textbf{Bước 2. Câu (i), $K=1,A=0,B=-1,f=x$.}
\[c_n=\frac2\pi\int_0^\pi x\sin nx\,dx
=\frac2\pi\left[-\frac{x\cos nx}n+\frac{\sin nx}{n^2}\right]_0^\pi
=\frac{2(-1)^{n+1}}n.\]
Vì vậy $u=\sum_{n\ge1}2(-1)^{n+1}e^{-(n^2+1)t}\sin nx/n$.

\textbf{Bước 3. Câu (ii), $K=1,A=-1,B=0,f=1$.}
$h=-1/2$, nên dữ liệu của $v$ là $e^{-x/2}$. Nguyên hàm $\int e^{hx}\sin nx=e^{hx}(h\sin nx-n\cos nx)/(h^2+n^2)$ cho
\[c_n=\frac2\pi\left[\frac{e^{-x/2}(-\sin nx/2-n\cos nx)}{n^2+1/4}\right]_0^\pi
=\frac{2n(1-(-1)^ne^{-\pi/2})}{\pi(n^2+1/4)}.\]
Suy $u=e^{x/2}\sum c_ne^{-(n^2+1/4)t}\sin nx$.

\textbf{Bước 4. Duy nhất và kiểm tra.}
Với hiệu $w$ có dữ liệu đầu $0$, biên $0$,
\begin{align*}
I'&=K[ww_x]_0^\pi-K\int w_x^2+\frac A2[w^2]_0^\pi+B\int w^2\\
&=-K\int w_x^2+2BI\le0
\end{align*}
trong cả hai câu vì $B\le0$. $I(0)=0$ suy $w=0$. Phép đổi ẩn và các ODE mode kiểm tra PDE; sine kiểm tra biên; Fourier trả lại dữ liệu trong miền. $f(\pi)=\pi$ ở (i), và $f(0)=f(\pi)=1$ ở (ii) không khớp biên $0$, nên không liên tục tại các góc này. Hai nghiệm trơn khi $t>0$ và đều tiến về $0$.

%ID VD_4_5_2_02
\textbf{Bước 1. Năng lượng bậc ba và tính duy nhất.}
Với $K>0$, $u_t=Ku_{xxx}$, $u(0)=u_x(0)=u(L)=0$, hai lần tích phân từng phần cho
\[I'=K\int uu_{xxx}
=K[uu_{xx}]_0^L-\frac K2[u_x^2]_0^L
=-\frac K2u_x(L,t)^2\le0.\]
Bản in bỏ hệ số $1/2$ ở đẳng thức cuối; dấu suy giảm vẫn đúng. Áp dụng cho hiệu hai nghiệm có mặt đầu $0$, suy $I=0$, tức duy nhất trong lớp nghiệm đủ trơn/năng lượng thích hợp.

\textbf{Bước 2. Vì sao không dùng cơ sở trực giao cấp hai.}
Tách biến $u=XT$ cho $X'''=\lambda X$, $T'=K\lambda T$ và ba biên $X(0)=X'(0)=X(L)=0$. Với hàm riêng phức, lấy phần thực của tích phân với $\overline X$:
\[\operatorname{Re}\lambda\int_0^L|X|^2dx=-\frac12|X'(L)|^2\le0.\]
Điều này không chứng minh $\lambda$ thực. Toán tử $D^3$ không tự liên hợp; liên hợp là $-D^3$ với biên $Y(0)=Y(L)=Y'(L)=0$, suy từ hạng $[X''\overline Y-X'\overline{Y'}+X\overline{Y''}]_0^L$. Vì vậy chưa được khẳng định dãy hàm riêng thực trong sách là cơ sở đầy đủ.
Nếu chỉ xét nhánh thực $\lambda=-k^3$, nghiệm thỏa hai biên trái là
\[X=e^{-kx}-e^{kx/2}\cos(\sqrt3kx/2)+\sqrt3e^{kx/2}\sin(\sqrt3kx/2)
=e^{-kx}+2e^{kx/2}\sin(\sqrt3kx/2-\pi/6).\]
Biên phải cho $e^{-3kL/2}=-2\sin(\sqrt3kL/2-\pi/6)$. Dấu trước sine trong $X$ phải là $+$; dấu $-$ ở công thức hàm riêng của bản in làm $X(0)=2$, sai biên.

\textbf{Bước 3. Giải đầy đủ bằng Laplace, không giả định cơ sở chưa chứng minh.}
Đặt $U(x,p)=\int_0^\infty e^{-pt}u(x,t)dt$, $\operatorname{Re}p>0$. Từ $pU-f=KU'''$:
\[U'''-\frac pK U=-\frac fK,\quad U(0,p)=U_x(0,p)=U(L,p)=0.\]
Đặt nghiệm cơ bản
\[S_p(x)=\sum_{m=0}^\infty\frac{(p/K)^m x^{3m+2}}{(3m+2)!}.\]
Đạo hàm chuỗi cho $S_p'''=(p/K)S_p$, $S_p(0)=S_p'(0)=0,S_p''(0)=1$. Biến thiên hằng số cho
\[U(x,p)=C(p)S_p(x)-\frac1K\int_0^xS_p(x-y)f(y)dy.\]
Đạo hàm tích phân ba lần tạo đúng $f(x)$ nhờ $S_p''(0)=1$. Biên $U(L)=0$ xác định
\[C(p)=\frac{\int_0^LS_p(L-y)f(y)dy}{KS_p(L)}.\]
Mẫu không bằng $0$ khi $\operatorname{Re}p>0$: nếu $S_p(L)=0$, $S_p$ là hàm riêng không tầm thường với trị riêng $p/K$ có phần thực dương, trái đẳng thức năng lượng.

\textbf{Bước 4. Công thức nghiệm và kiểm tra.}
\[
u(x,t)=\mathcal L^{-1}_{p\to t}
\left\{\frac{S_p(x)}{KS_p(L)}\int_0^LS_p(L-y)f(y)dy
-\frac1K\int_0^xS_p(x-y)f(y)dy\right\}.
\]
Ở đây $\mathcal L^{-1}H=(2\pi i)^{-1}\int_{\gamma-i\infty}^{\gamma+i\infty}e^{pt}H(p)dp$, $\gamma>0$, hiểu theo phép đảo Laplace trong lớp năng lượng; với dữ liệu trơn tương thích có nghiệm cổ điển. Công thức $U$ thỏa ODE, cả ba biên và hạng mặt đầu $f$ trong $pU-f=KU'''$, nên phép đảo trả lại đúng bài toán. Với $f\in L^2$, nghiệm được hiểu là nửa nhóm co sinh bởi $KD^3$ trên miền $H^3$ có ba biên trên; tiêu tán vừa chứng minh và nghiệm resolvent cho mọi $p>0$ bảo đảm tồn tại và duy nhất. Đây là công thức giải cho dữ liệu tổng quát mà không bỏ các thành phần phổ phức hoặc giả định các hệ số trực giao.

%ID VD_4_6_2_01
\textbf{Bước 1. Nhân nhiệt trên toàn trục.}
Đặt $g_K(z,t)=(4\pi Kt)^{-1/2}e^{-z^2/(4Kt)}$. Nghiệm $u=\int_{\mathbb R}g_K(x-y,t)f(y)dy$. Tích phân của $g_K$ bằng $1$, và biến đổi Fourier của nó là $e^{-Kt\xi^2}$, nên một mode $\cos ky$ biến thành $e^{-Kk^2t}\cos kx$.
\textbf{Bước 2. Câu (i), $K=1,f=\cos^23x$.}
$\cos^23x=(1+\cos6x)/2$, suy
\[u=\frac12+\frac12e^{-36t}\cos6x.\]
Kiểm tra $u_t=-18e^{-36t}\cos6x=u_{xx}$, $u(x,0)=\cos^23x$.
\textbf{Bước 3. Câu (ii), $K=2,f=e^{-x^2}$.}
\[u=\frac1{\sqrt{8\pi t}}\int_{\mathbb R}
\exp\left(-\frac{(x-y)^2}{8t}-y^2\right)dy.\]
Hoàn thành bình phương:
\[-\frac{(x-y)^2}{8t}-y^2
=-\frac{1+8t}{8t}\left(y-\frac{x}{1+8t}\right)^2-\frac{x^2}{1+8t}.\]
Đặt $z=\sqrt{(1+8t)/(8t)}(y-x/(1+8t))$, $dy=\sqrt{8t/(1+8t)}dz$. Tích phân $\int e^{-z^2}dz=\sqrt\pi$ cho
\[u=\frac1{\sqrt{1+8t}}\exp\left(-\frac{x^2}{1+8t}\right).\]
\textbf{Bước 4. Kiểm tra đạo hàm.} Đặt $D=1+8t$. $u_t/u=-4/D+8x^2/D^2$, $u_{xx}/u=-2/D+4x^2/D^2$, nên $u_t=2u_{xx}$. Khi $t\downarrow0$, $D\to1$ trả lại $f$.

%ID VD_4_6_2_02
\textbf{Bước 1. Thác triển lẻ.}
Với $u(0,t)=0$, đặt $\widetilde f(y)=f(y)$ khi $y>0$, $\widetilde f(y)=-f(-y)$ khi $y<0$. Vì Gauss chẵn, nghiệm toàn trục của dữ liệu lẻ vẫn lẻ, nên bằng $0$ tại $x=0$. Tách tích phân và đổi $y\mapsto-y$ trên nửa âm:
\[u=\int_0^\infty[g_K(x-y,t)-g_K(x+y,t)]f(y)dy.\]
\textbf{Bước 2. Câu cụ thể $K=2,f=e^{-y}$.}
Đặt $I_\pm=\int_0^\infty g_2(x\mp y,t)e^{-y}dy$, trong đó $I_+$ dùng $x-y$, $I_-$ dùng $x+y$. Hoàn thành bình phương:
\begin{align*}
-\frac{(x-y)^2}{8t}-y&=-\frac{(y+4t-x)^2}{8t}+2t-x,\\
-\frac{(x+y)^2}{8t}-y&=-\frac{(y+4t+x)^2}{8t}+2t+x.
\end{align*}
Đổi biến $z=(y+4t\mp x)/\sqrt{8t}$; cận $y=0$ trở thành $(4t\mp x)/\sqrt{8t}$, và dùng $\int_a^\infty e^{-z^2}dz=\sqrt\pi\,\operatorname{erfc}(a)/2$:
\[I_+=\tfrac12e^{2t-x}\operatorname{erfc}\frac{4t-x}{\sqrt{8t}},\quad
I_-=\tfrac12e^{2t+x}\operatorname{erfc}\frac{4t+x}{\sqrt{8t}}.\]
Vậy
\[u=\frac12e^{2t-x}\operatorname{erfc}\frac{4t-x}{\sqrt{8t}}
-\frac12e^{2t+x}\operatorname{erfc}\frac{4t+x}{\sqrt{8t}}.\]
\textbf{Bước 3. Kiểm tra.} Tại $x=0$ hai hạng triệt tiêu. Với $x>0,t\downarrow0$, đối số thứ nhất tiến về $-\infty$ nên erfc tiến về $2$, đối số thứ hai về $\infty$ nên erfc về $0$; suy $u\to e^{-x}$. Nhân ảnh giải PDE nên $u_t=2u_{xx}$. $f(0)=1$ không khớp biên $0$, nên không liên tục ở góc $(0,0)$; nghiệm trơn với $x\ge0,t>0$.

%ID VD_4_6_2_03
\textbf{Bước 1. Chuyển bài $u_t=u_{xx}+xu$ về nhiệt.}
Đặt $u=e^{tx+t^3/3}v(z,t)$, $z=x+t^2$. Tính
\begin{align*}
u_t&=e^{tx+t^3/3}\{(x+t^2)v+2tv_z+v_t\},\\
u_{xx}+xu&=e^{tx+t^3/3}\{(t^2+x)v+2tv_z+v_{zz}\}.
\end{align*}
Do đó $v_t=v_{zz}$, $v(z,0)=f(z)$. Phép đổi biến này biện minh công thức cho $f$ liên tục bị chặn mà không cần đánh giá $\widehat f$ tại tần số phức.
\textbf{Bước 2. Fourier cho bài đã đổi.} Với quy ước $\widehat v=\int e^{-i\xi z}v\,dz$, $\widehat v_t=-\xi^2\widehat v$, suy $\widehat v=e^{-t\xi^2}\widehat f$. Biến đổi ngược cho $v=g_1*f$. Vì vậy
\[u=e^{tx+t^3/3}\int_{\mathbb R}g_1(x+t^2-y,t)f(y)dy.\]
Gom mũ:
\[tx+t^3/3-\frac{(x+t^2-y)^2}{4t}
=-\frac{(x-y)^2}{4t}+\frac{t(x+y)}2+\frac{t^3}{12}.\]
Vậy
\[u=\int_{\mathbb R}G(x,y,t)f(y)dy,\qquad
G=\frac1{\sqrt{4\pi t}}
\exp\left[-\frac{(x-y)^2}{4t}+\frac{t(x+y)}2+\frac{t^3}{12}\right].\]
\textbf{Bước 3. Kiểm tra Green.} $G>0$, đối xứng theo $x,y$, trơn với $t>0$. Với $\ell=\log G$,
\[\ell_x=-\frac{x-y}{2t}+\frac t2,\quad \ell_{xx}=-\frac1{2t},\quad
\ell_t=-\frac1{2t}+\frac{(x-y)^2}{4t^2}+\frac{x+y}2+\frac{t^2}4.\]
Ta có $\ell_t=\ell_{xx}+\ell_x^2+x$, suy $G_t=G_{xx}+xG$. Từ biểu diễn Gauss đã dịch,
\[\int_{\mathbb R}G(x,y,t)dy=e^{tx+t^3/3}.\]
Nhân không có tổng bằng $1$ do hạng phản ứng $xu$.
\textbf{Bước 4. Dữ liệu đầu và sự tập trung.} Với $x$ trong tập compact, hệ số $e^{tx+t^3/3}\to1$ và tâm $x+t^2\to x$. Gauss là xấp xỉ đơn vị, nên $u(x,t)\to f(x)$ tại điểm liên tục, cũng liên tục khi $(x,t)\to(x_0,0)$. Với $\delta>0$, tích phân nhân trên $|y-x|>\delta$ tiến về $0$. Điều kiện này phải là khoảng cách $|y-x|$, không phải $|x|$ như một dòng sai biến trong bản in. Công thức không bảo đảm nghiệm bị chặn theo $x$ khi $t>0$, vì phản ứng tuyến tính theo $x$ có thể làm nghiệm tăng theo mũ.

%ID VD_4_6_2_04
\textbf{Bước 1. Nhận dạng một mode cộng hưởng.}
Nguồn $e^{-t}\cos x$ và mặt đầu $0$ cho phép đặt $u=T(t)\cos x$. Khi đó $u_{xx}=-T\cos x$, PDE thành $T'+T=e^{-t}$, $T(0)=0$.
\textbf{Bước 2. Thực hiện tích phân.} Nhân với $e^t$:
\[(e^tT)'=1,\quad e^tT=t,\quad T=te^{-t}.\]
Tương đương Duhamel: $u=\int_0^te^{-(t-s)}e^{-s}\cos x\,ds=e^{-t}\cos x\int_0^tds$.
\textbf{Bước 3. Kiểm tra.}
\[u=te^{-t}\cos x,\quad u_t=(1-t)e^{-t}\cos x,\quad
u_{xx}=-te^{-t}\cos x.\]
Suy $u_t-u_{xx}=e^{-t}\cos x$ và $u(x,0)=0$. Thừa số $t$ xuất hiện vì tốc độ giảm của nguồn bằng tốc độ giảm của mode nhiệt.

%ID VD_4_6_2_05
\textbf{(a) Nhân nhiều chiều tách tích.}
Với $u_t=K\Delta u$ trên $\mathbb R^n$, nhân là
\[g_K^{(n)}(x-y,t)=\prod_{j=1}^n g_K(x_j-y_j,t).\]
Dữ liệu $f(y)=\prod f_j(y_j)$ cho, nhờ Fubini,
\[u(x,t)=\int_{\mathbb R^n}\prod_{j=1}^n[g_K(x_j-y_j,t)f_j(y_j)]dy
=\prod_{j=1}^nu_j(x_j,t).\]
Kiểm tra trực tiếp cũng được: $u_t=\sum_j(u_j)_t\prod_{i\ne j}u_i$, còn $K\Delta u=\sum_jK(u_j)_{x_jx_j}\prod_{i\ne j}u_i$, bằng nhau vì mỗi $u_j$ giải nhiệt một chiều. Tại $t=0$, tích trả lại $f$.
\textbf{(b) $n=2,K=1,f_1=\cos x_1,f_2=\sin x_2$.}
Mode một chiều cho $u_1=e^{-t}\cos x_1$, $u_2=e^{-t}\sin x_2$. Vì vậy
\[u=e^{-2t}\cos x_1\sin x_2.\]
$u_t=-2u$, $\Delta u=-u-u=-2u$, và mặt đầu đúng. Tốc độ giảm tổng là tổng hai trị riêng $1+1=2$.
